\section{Discussion}
\label{sec:discussion}

Table~\ref{tab:verdicts} sets the results of E1--E5 against the hypotheses of
Section~\ref{sec:definition}. None is supported as stated. Four are supported in part, mostly
because the predicted effect appeared for one kind of sub-task or one indicator and not for the
others, and H3 is not supported.

\begin{table}[t]
  \centering
  \caption{Verdicts on the hypotheses. Partly: some of the predictions hold and others do not or
  were not tested. Evidence from Section~\ref{sec:experiments}; intervals are 95\% task-cluster
  bootstrap intervals.}
  \label{tab:verdicts}
  \small
  \begin{tabularx}{\textwidth}{@{}l l X@{}}
    \toprule
    Hypothesis & Verdict & Evidence \\
    \midrule
    H1 dose--response & Partly & Pass rate falls with $k$ for cascades in both models and for four
      user requests in \texttt{gpt-oss}, beyond the filler control for \texttt{qwen3-coder}
      cascades; flat for \textsf{env}, \textsf{text} and filler. The off-goal share rises with $k$
      only for user requests (E1). \\
    H2 source signatures & Partly & User requests cause incompletion (\texttt{gpt-oss}, $k=4$), as
      predicted. \textsf{Text} causes no scope violations, \textsf{env} no edits to protected tests,
      and cascades are the most damaging source, not the least (E1). \\
    H3 mitigations & Not supported & No mitigation changes the pass rate after Holm correction;
      deterministic hooks do no better than instruction-based mitigations (E2). \\
    H4 compaction & Partly & Forced compaction raises the post-compaction off-goal share by $+0.13$
      $[+0.06, +0.20]$, with an imperfect placebo; the pass rate changes by $-0.05$
      $[-0.11, +0.00]$; no constraint violations; recall not measurable; inheritance not tested
      (E3 pilot). \\
    H5 observability & Partly & Summaries over-claim in 42--50\% of baseline episodes and up to 81\%
      under cascades, and usually omit the out-of-scope files they changed. A monitor flags more
      failing episodes per review than review at every pause, but only 18--19\% of them early;
      monitor degradation not tested (E4). \\
    \bottomrule
  \end{tabularx}
\end{table}

\subsection{What kind of drift did we observe?}

The literature that motivated this study frames goal drift as the gradual \emph{displacement} of
the assigned goal by something else: competing objectives, environmental pressure, inherited
context \citep{arike2025goaldrift,menon2026inherited,saebo2026asymmetric}. Our results support a
narrower and less dramatic picture, with three features.

\paragraph{The kind of sub-task matters more than how many there are.} Up to four injected
sub-tasks of most kinds left both models where they started. Neither neutral interruptions of equal
length nor distractors moved the pass rate. Only two kinds of sub-task hurt: dependencies that had to
be repaired before the goal could be reached, and, for the weaker model, a stack of four sanctioned
side requests. A count of sub-tasks is therefore a poor predictor of drift. What matters is whether
a sub-task demands real work on the path to the goal or competes with the goal for the agent's
limited budget of turns.

\paragraph{The loss appears as incompletion, not as deviation.} In no condition did the agents edit
existing files outside the admissible scope to any meaningful degree (below 3.5\% of changed lines).
They did not chase planted \texttt{TODO}s into unrelated modules, nor did they ``fix'' the unrelated
failing tests at the expense of the goal. The goal was lost by \emph{not finishing it}. With four
side requests, \texttt{gpt-oss} completed fewer requirements of the goal and fewer of the requests.
With broken dependencies, both models more often ended with requirements failing. This matches the
\emph{unreliability} account of multi-turn degradation \citep{laban2025lost} better than a
goal-replacement account. The judge panel of E5 points the same way
(Section~\ref{sec:judge-results}). The work it labels as leaving the goal grows with the number of
side requests, to 0.53 (\texttt{qwen3-coder}) and 0.62 (\texttt{gpt-oss}) of the segments at
$k=4$, but almost all of it is work the user
asked for or repairs the goal needed. Work on distractors stays below 7\% of segments in every
cell. In these runs, then, drift is mostly legitimate widening of the scope at the user's request,
followed by goals left unfinished, and not the pursuit of something illegitimate. Two limits apply.
The panel's labels are reliable, not validated against humans, and the judges could not agree on
goal recall from the probe, so whether the agents' \emph{beliefs} about the goal shifted remains
open; the exploratory goal-replacement label shows no shift either.

\paragraph{Necessary work is the most dangerous.} H2 predicted that necessary cascades would cause
the least drift, because resolving them serves the goal. They caused the most. Two explanations are
compatible with the data:
\begin{itemize}
  \item the agent never found or fixed the broken dependency, so the loss is unresolved necessary
        work: a capability failure that the injection made visible;
  \item the agent fixed it but spent its budget, or its attention, doing so, and did not return to
        the remaining requirements: drift in the intended sense.
\end{itemize}
The two predict different trajectories: the first ends with the dependency still broken, the second
with it repaired and requirements missing. The judge labels separate them only in part. For
\texttt{gpt-oss} the share of segments spent off the goal does not rise with cascades, which fits
the first explanation: little visible repair work. For \texttt{qwen3-coder} it rises to 18\% at
$k=4$, none of it labelled as distraction, which leaves room for the second. Whether each
dependency ended up repaired was not judged. Either way the practical lesson is the same: the
sub-tasks that look most
legitimate, and that a supervisor is least likely to question, are where goal completion is most at
risk.

The resistance to distractors deserves a caveat. Studies of benign distraction and of instruction
hijacking report strong effects
\citep{ma2024caution,mavali2026tab,debenedetti2024agentdojo}. Our distractors were mild and
plausible rather than adversarial, they pointed away from the task, and every goal ended with ``keep
the change focused on this feature''. Under these conditions both models behaved conservatively.
The result says little about more salient or adversarially phrased material.

\paragraph{Compaction shifts effort from the goal to the requests.} The pilot of
Section~\ref{sec:e3} suggests a mechanism behind part of the incompletion. After the context was
summarised, \texttt{qwen3-coder} gave more of its remaining segments to the user's side
requests and slightly less to the goal, with a small, uncertain loss on the goal and a small gain on
the requests. This fits the concern of Section~\ref{sec:drift} that summarisation decides what
survives of the original request. One reading, which we have not checked against the compaction
summaries themselves, is that a summary written mid-session keeps the recent, concrete requests and
gives less weight to the original specification. The extra work is still sanctioned, so a
supervisor would see nothing wrong in it. The observational data cannot confirm the effect, because spontaneous compaction marks episodes that
were already in trouble, and the pilot is small, with a placebo that is not clean. The effect needs a
larger, preregistered test before it counts as established.

\subsection{Implications for the supervisor}

\paragraph{Drift of this kind is quiet.} A supervisor who reviews the diff---the most common form
of oversight (Section~\ref{sec:hai})---would see a focused, in-scope change in almost every episode,
including those in which the goal was badly incomplete. The failure is visible only against the
requirements, not against the code. This is the situation human-factors research warns about: a
supervisor's situation awareness depends on the signals the system exposes \citep{endsley1995situation},
and a tidy diff with a confident summary invites trust that is not calibrated to the outcome
\citep{lee2004trust}. The summaries add to the problem, which bears on RQ4 and H5. In E1, 42\%
of \texttt{qwen3-coder} baseline episodes and 50\% of \texttt{gpt-oss} ones end with a summary
that claims completion while a hidden test fails. With cascades the rate rises to about 0.75--0.81
in both models, and in E2 the goal re-anchoring hook raises it from 0.69 to 0.89. Part of the
cascade increase is mechanical, since more episodes fail a test, but it means that the sub-tasks
most likely to cost the goal are also those after which the agent's report is least reliable. These
intervals are not corrected for multiplicity, and over-claiming is measured at the level of the
episode, not per requirement. Relying on passing tests does not close the gap. In 19\% of E1 episodes the
entire visible suite passed while some requirements of the goal failed, and agents are known to optimise for the tests they can see
\citep{zhong2025impossiblebench}. Oversight tools should therefore show progress \emph{per
requirement}, not per file, and make what is still unverified explicit. This is a concrete form of
the ``conveying progress'' challenge of \citet{bansal2024communication}.

\paragraph{What a supervisor can catch (E4).} The supervisor simulation (Section~\ref{sec:e4})
puts numbers on this argument. Even with segment labels applied without error, a supervisor who
reviews what the agent is doing catches few failing episodes in time. Reviewing after every segment
flags about half of them before the end, at six reviews per episode and with an alert in most
passing episodes as well, because the alerts fire on detours the user asked for. A monitor for
illegitimate work is cheap but flags fewer than one in five. Reading the final summary flags fewer
than one in five at the end, since the summary claims completion in most failing episodes, and it
usually does not name the out-of-scope files that were changed. Because these failures are
incompletion, the activity looks normal. The check that would separate failing from passing
episodes is the one the agent's report replaces: running the acceptance tests of each requirement.
A supervisor's interface should show that check, and the harness should run it before the agent is
allowed to report completion.

\paragraph{Where to put checkpoints.} The losses concentrate at recognisable moments: when the agent
starts repairing something the goal depends on, and when side requests pile up. These are natural
triggers for the mixed-initiative interruptions advocated by \citet{bui2026proactivity}
\citep[see also][]{horvitz1999mixed}. They also suit a dynamic autonomy policy of the kind
implemented by Hedwig \citep{shukla2026hedwig}, which tightens oversight in unfamiliar territory. A
harness can detect both moments cheaply: an edit to a file outside $S(\Gz)$ that makes a failing
import or test pass, or a third pending user request. It can then ask the human whether to continue,
defer or delegate. In E4, review at every pause spent six reviews per episode and alerted in 55--77\% of
passing episodes, so periodic review mostly interrupts work that was going well. The triggers
proposed here are untested; the E4 monitor was triggered by illegitimate work, which is rare in
these runs.

\paragraph{The user is a source of drift too.} Sanctioned side requests are the user's own doing.
For the weaker model, four of them in one session cost half of the goal's completion and more than
half of the requests themselves. A design response is to treat a mid-session request as a new item
in a queue rather than an interruption of the current one: completed after the goal, or handed to a
separate session with its own context. Claude Code already has the ingredients, namely task lists,
background sessions and sub-agents. Our results suggest that these work only if the harness, not the
model, decides to use them.

\paragraph{Mitigations should not depend on compliance.} E2 does not support H3 in its strong form.
Deterministic hooks were not better than instruction-based mitigations, but the reasons differ.
\begin{itemize}
  \item \emph{Re-anchoring} was applied reliably, yet it repeats information the model already has
        and does nothing about the work that competes with the goal.
  \item The \emph{stop-hook checklist} acts only when the agent decides to stop. Under turn limits
        that rarely happened, and in real use agents also stop for other reasons (budgets,
        timeouts, errors).
  \item \emph{Delegation} would isolate sub-tasks as intended \citep{fei2026codedelegator}, but the
        model delegated in two episodes out of 36.
  \item The \emph{task list} was adopted more often and gave the largest, still non-significant,
        gain, at the cost of more tool calls and more stray files.
\end{itemize}
The common thread is that lightweight interventions inside the session do not change \emph{what}
the agent spends its budget on. More promising directions are interventions with teeth:
\begin{itemize}
  \item a completion gate that checks requirements, not just a stop;
  \item harness-enforced delegation of recognised sub-task types, rather than a suggestion to
        delegate;
  \item budgeting that reserves turns for the goal.
\end{itemize}
In all three the compliance of the model stops being a free variable.

\subsection{Limitations}

\begin{description}[leftmargin=0pt]
  \item[Models.] The agents are 20--30B open-weight models, not the Claude models that Claude Code
        is designed for. The harness, system prompt and tool descriptions are tuned for those, and the
        local models often used up their turns: 14\% of \texttt{qwen3-coder}'s final segments ended
        at the turn limit. Stronger models may absorb more sub-tasks
        before losing the goal, which would shift the thresholds, not necessarily remove them. The
        two models also differ widely in baseline competence, so cross-model comparisons are about
        patterns, not magnitudes.
  \item[Benchmark.] DriftBench-mini is small and synthetic: twelve libraries of a few hundred
        lines, with goals specified as numbered requirements and sub-tasks written by us. Real
        sessions are longer, goals are vaguer and sub-tasks arise organically. Some cascades are
        necessary only for some implementations, one is all-or-nothing, and two can never change the
        score (Section~\ref{sec:bench}). Because cascades widen the scope by design, part of their
        effect may reflect task difficulty rather than drift; the judge-based measures
        (Section~\ref{sec:judge-results}) address this only in part.
  \item[Protocol.] The fixed structure of pauses and per-segment turn limits makes conditions
        comparable but is not how people use agents. It also shaped the mitigations: the stop hook,
        for instance, rarely had a stop to act on. Results may differ in free-running sessions with
        larger budgets.
  \item[Judges.] The judge-based measures rest on a panel of three LLM judges. Their agreement
        shows reliability, not validity: no human labelled the trajectories. The binary off-goal
        label was chosen as the primary measure after seeing its held-out agreement, and codebook v2
        was written after reading the disagreements on the full set. Reliability is lower on
        \texttt{qwen3-coder} trajectories ($\kappa=0.615$ for the binary label on the full set, 0.48 for
        the five labels). The judges cannot be blinded to the condition, because they need the injected
        sub-tasks to tell sanctioned from distracting work. In 3 of 1{,}332 episodes (0.2\%) the
        agent read the harness's log files with \texttt{grep -R ..} from its sandbox. None of
        them read the episode specification that names the condition, and agent model names that
        leaked into two judge items were scrubbed. Goal recall is not used ($\kappa=0.17$) and goal
        replacement is exploratory ($\kappa=0.45$).
  \item[Measures and power.] Goal recall is not reported and goal replacement only as exploratory
        (see above); report fidelity is measured per episode, not per requirement. With 36 episodes per
        cell, the confidence intervals on a difference are about $\pm 0.05$--$0.10$ for
        \texttt{qwen3-coder} and up to $\pm 0.2$ for the more variable \texttt{gpt-oss}, so moderate effects, including the E2
        mitigations, cannot be excluded.
  \item[E3.] The forced-compaction pilot has one model, one source and twelve tasks with 36
        episodes per arm. Its placebo comparison is not clean ($+0.05$ on segments that precede the
        compaction), so part of the post-compaction difference may be imbalance between arms.
        Agreement on \texttt{claims\_complete} in E3 ($\kappa=0.56$) is below the criterion, so its
        over-claim figures are weaker than those of E1 and E2. The observational comparison is
        correlational, since compaction occurs in long, struggling episodes. Inherited drift was
        not tested.
  \item[E4.] The supervisor is simulated, not observed. It applies the panel labels without error,
        so the figures are upper bounds for each policy; real reviewers misread transcripts and tire.
        The P1 alert rule counts sanctioned and necessary detours as alerts, which inflates its false
        alarms, and omitted changes rest on a strict file-name match, which inflates omissions.
        Whether monitors degrade on longer transcripts was not tested.
  \item[Infrastructure.] The last 45 E1 episodes ran on a different cluster with identical model
        weights but without container isolation. They are a small share of one model's runs, and
        no episode referenced evaluation material, but the change is a deviation from the protocol.
\end{description}
